% Test of the caption package
% (c) 2003-2004 Axel Sommerfeldt

\documentclass[a4paper]{article}
\setlength\parindent{0pt}
\setlength\parskip{\medskipamount}

\errorcontextlines9
\def\StartTracing{\typeout{StartTracing}\tracingmacros=1\tracingcommands=1\relax}
\def\StopTracing{\tracingmacros=0\tracingcommands=0\relax\typeout{StopTracing}}

\DeclareRobustCommand\cs[1]{\texttt{\char`\\\string#1}}
%\providecommand\star{\ttfamily*}
%\providecommand\marg[1]{{\ttfamily\char`\{#1\char`\}}}
%\providecommand\oarg[1]{{\ttfamily[#1]}}
%\providecommand\parg[1]{{\ttfamily(#1)}}
\def\ifundefined#1{\expandafter\ifx\csname#1\endcsname\relax}

\usepackage{shortvrb}\MakeShortVerb{\|}
%\usepackage{pkgindoc}
\makeatletter
\let\@onlypreamble\@gobble
\let\Orig@makecaption\@makecaption
\makeatother

\usepackage[bf,sl,Large,hang,justification=raggedright]{caption}[2003/12/14]
\DeclareCaptionOption{debug}{}
\DeclareCaptionFormat{normal}{#1#2#3\par}
%\captionsetup{margin=\leftmargini}
%
%\usepackage{subfig}
%
\usepackage{float}%row}
\usepackage{listings}[2004/02/11]
\usepackage{longtable}\makeatletter\let\Orig@LT@makecaption\LT@makecaption\makeatother
\usepackage{rotating}
\usepackage{sidecap}
\usepackage{supertabular}\makeatletter\let\Orig@ST@caption\ST@caption\makeatother

\ifundefined{floatstyle}
\else
  \floatstyle{ruled}
  \newfloat{Program}{tbp}{lop}[section]
  \floatstyle{plain}
  \newfloat{Example}{t}{lox}[section]
\fi

\long\def\cap#1{%
  \shortcap{#1}\longcap{#1}}
\long\def\shortcap#1{%
  \bgroup
  #1%
  \captionof{figure}[]{\shorttext}%
  \egroup}
\long\def\longcap#1{%
  \bgroup
  #1%
  \captionof{figure}[]{\longtext}%
  \egroup}
\def\shorttext{Short caption}
\def\longtext{Long long long long long long long long
  long long long long long long long long long long long long
  long long long long long long long long long long long caption}

\newif\iftesterrors
%\testerrorstrue

\let\oldsection\section
\def\section#1{\typeout{#1}\oldsection{#1}}

\begin{document}

\title{Test of the caption package}
\author{Axel Sommerfeldt}
\maketitle


%************************************%
%*                                  *%
%* Options                          *%
%*                                  *%
%************************************%

\section{Options}
|\usepackage[bf,sl,Large,hang,justification=raggedright]{caption}|

\noindent Normal Text
\cap{}
\captionsetup{style=default}


%************************************%
%*                                  *%
%* Obsolete stuff                   *%
%*                                  *%
%************************************%

\clearpage
\section{Obsolete stuff}

|\renewcommand\captionsize{\small\itshape}|
\cap{\renewcommand\captionsize{\small\itshape}}

|\renewcommand\captionfont{\small\itshape}|
\cap{\renewcommand\captionfont{\small\itshape}}

|\renewcommand\captionlabelfont{\bfseries}|
\cap{\renewcommand\captionlabelfont{\bfseries}}

|\renewcommand\captiontextfont{\itshape}|
\cap{\renewcommand\captiontextfont{\itshape}}

|\setlength\captionmargin{1in}|
\cap{\setlength\captionmargin{1in}}


%************************************%
%*                                  *%
%* Vertical spacing and positioning *%
%*                                  *%
%************************************%

\clearpage
\section{Vertical spacing and positioning}

|\captionsetup[figure]{position=b}|\\
|\captionsetup[table]{position=t}|
\captionsetup[figure]{position=b}
\captionsetup[table]{position=t}

\vbox{\hrule
\captionof{figure}[]{Short caption}%
\hrule
\captionof{table}[]{Short caption}%
\hrule}

|\captionsetup{aboveskip=20pt,belowskip=10pt}|
\captionsetup{aboveskip=20pt,belowskip=10pt}

\vbox{\hrule
\captionof{figure}[]{Short caption}%
\hrule
\captionof{table}[]{Short caption}%
\hrule}

|\captionsetup{aboveskip=10pt,belowskip=0pt,position=auto,debug=3}|\\
|\showcaptionsetup{table}|\\
|\clearcaptionsetup{table}|\\
\captionsetup{aboveskip=10pt,belowskip=0pt,position=auto,debug=3}
\showcaptionsetup{table}

\clearcaptionsetup{figure}
\clearcaptionsetup{table}

\begin{figure}[!ht]
\centering Text
\caption{This is a figure}
\end{figure}

%\begin{figure}[!ht]
%\centering Text\par
%\caption{This is a figure}
%\end{figure}

\begin{table}[!ht]
\caption{This is a table}
\centering\begin{tabular}{ll}
A & B \\
C & D \\
\end{tabular}
\end{table}

\captionsetup{style=default,position=default,aboveskip=10pt,belowskip=0pt,debug=1}


%************************************%
%*                                  *%
%* Styles                           *%
%*                                  *%
%************************************%

\clearpage
\section{Styles}

|\DeclareCaptionStyle{italic}%|\\
|  {labelfont={sf,bf},textfont={rm,it},indention=18pt,%|\\
|   labelsep=period,justification=raggedright}|\\
|\captionsetup{style=italic}|

\DeclareCaptionStyle{italic}%
  {labelfont={sf,bf},textfont={rm,it},indention=18pt,%
   labelsep=period,justification=raggedright}
\cap{\captionsetup{style=italic}}

|\DeclareCaptionStyle{italix}[justification=raggedleft]%|\\
|  {labelfont={sf,bf},textfont={rm,it},indention=18pt,%|\\
|   labelsep=period,justification=raggedright}|\\
|\captionsetup{style=italix}|

\DeclareCaptionStyle{italix}[justification=raggedleft]%
  {labelfont={sf,bf},textfont={rm,it},indention=18pt,%
   labelsep=period,justification=raggedright}
\cap{\captionsetup{style=italix}}

|\DeclareCaptionStyle{nooneline}{}|\\
|\captionsetup{style=nooneline}|

\DeclareCaptionStyle{nooneline}{}
\cap{\captionsetup{style=nooneline}}

|\captionsetup{style=default}|

\cap{\captionsetup{style=default}}


%************************************%
%*                                  *%
%* Formats                          *%
%*                                  *%
%************************************%

\clearpage
\section{Formats}

|\captionsetup{format=default}|
\cap{\captionsetup{format=default}}

|\captionsetup{format=default,indention=\leftmargini}|
\cap{\captionsetup{format=default,indention=\leftmargini}}

|\captionsetup{format=hang}|
\cap{\captionsetup{format=hang}}

|\captionsetup{format=hang,indention=10pt}|
\cap{\captionsetup{format=hang,indention=10pt}}

|\captionsetup{format=hang,indention=-10pt}|
\cap{\captionsetup{format=hang,indention=-10pt}}

|\DeclareCaptionFormat{hfill}{#1\hskip 10pt plus 1fill\relax#3}|\\
|\captionsetup{format=hfill}|
\DeclareCaptionFormat{hfill}{#1\hskip 10pt plus 1fill\relax#3}
\cap{\captionsetup{format=hfill}}


%************************************%
%*                                  *%
%* Label formats                    *%
%*                                  *%
%************************************%

\clearpage
\section{Label formats}

|\captionsetup{labelformat=empty}|
\cap{\captionsetup{labelformat=empty}}

|\captionsetup{labelformat=simple}|
\cap{\captionsetup{labelformat=simple}}

|\captionsetup{labelformat=parens}|
\cap{\captionsetup{labelformat=parens}}

|\DeclareCaptionLabelFormat{mylabelformat}{#1\textvisiblespace#2}|\\
|\captionsetup{labelformat=mylabelformat}|

\DeclareCaptionLabelFormat{mylabelformat}{#1\textvisiblespace#2}
\cap{\captionsetup{labelformat=mylabelformat}}

|\def\figurename{}|

\cap{\def\figurename{}\captionsetup{labelformat=mylabelformat}}

|\DeclareCaptionLabelFormat{mylabelformat}{\bothIfFirst{#1}{\textvisiblespace}#2}|\\
|\captionsetup{labelformat=mylabelformat}|

\DeclareCaptionLabelFormat{mylabelformat}{\bothIfFirst{#1}{\textvisiblespace}#2}
\cap{\captionsetup{labelformat=mylabelformat}}

|\def\figurename{}|

\cap{\def\figurename{}\captionsetup{labelformat=mylabelformat}}


%************************************%
%*                                  *%
%* Label separators                 *%
%*                                  *%
%************************************%

\clearpage
\section{Label separators}

|\captionsetup{labelsep=none}|
\cap{\captionsetup{labelsep=none}}

%|\captionsetup{labelsep=default}|
%\cap{\captionsetup{labelsep=default}}

|\captionsetup{labelsep=colon}|
\cap{\captionsetup{labelsep=colon}}

|\captionsetup{labelsep=period}|
\cap{\captionsetup{labelsep=period}}

|\captionsetup{labelseparator=newline}|
\cap{\captionsetup{labelseparator=newline}}

%|\DeclareCaptionLabelSeparator{linebreak}{\\}|\\
%|\captionsetup{labelseparator=linebreak}|
%\DeclareCaptionLabelSeparator{linebreak}{\\}
%\cap{\captionsetup{labelseparator=linebreak}}

|\DeclareCaptionLabelSeparator{visiblespace}{\textvisiblespace}|\\
|\captionsetup{labelseparator=visiblespace}|
\DeclareCaptionLabelSeparator{visiblespace}{\textvisiblespace}
\cap{\captionsetup{labelseparator=visiblespace}}

|\DeclareCaptionLabelSeparator{en-dash}{\textrm{ -- }}|\\
|\captionsetup{labelseparator=en-dash}|
\DeclareCaptionLabelSeparator{en-dash}{\textrm{ -- }}
\cap{\captionsetup{labelseparator=en-dash}}


%************************************%
%*                                  *%
%* Justifications                   *%
%*                                  *%
%************************************%

\clearpage
\section{Justifications}
\captionsetup{singlelinecheck=0}

|\captionsetup{format=default,justification=justified}|
\cap{\captionsetup{format=default,justification=justified}}

|\captionsetup{format=default,justification=centering}|
\cap{\captionsetup{format=default,justification=centering}}

|\captionsetup{format=default,justification=centerfirst}|
\cap{\captionsetup{format=default,justification=centerfirst}}

|\captionsetup{format=default,justification=centerlast}|
\cap{\captionsetup{format=default,justification=centerlast}}

|\captionsetup{format=default,justification=raggedleft}|
\cap{\captionsetup{format=default,justification=raggedleft}}

|\captionsetup{format=default,justification=raggedright}|
\cap{\captionsetup{format=default,justification=raggedright}}

|\captionsetup{format=default,justification=Centering}|
\cap{\captionsetup{format=default,justification=Centering}}

|\captionsetup{format=default,justification=RaggedLeft}|
\cap{\captionsetup{format=default,justification=RaggedLeft}}

|\captionsetup{format=default,justification=RaggedRight}|
\cap{\captionsetup{format=default,justification=RaggedRight}}

|\captionsetup{format=hang,justification=justified}|
\cap{\captionsetup{format=hang,justification=justified}}

|\captionsetup{format=hang,justification=centering}|
\cap{\captionsetup{format=hang,justification=centering}}

|\captionsetup{format=hang,justification=centerfirst}|
\cap{\captionsetup{format=hang,justification=centerfirst}}

|\captionsetup{format=hang,justification=centerlast}|
\cap{\captionsetup{format=hang,justification=centerlast}}

|\captionsetup{format=hang,justification=raggedleft}|
\cap{\captionsetup{format=hang,justification=raggedleft}}

|\captionsetup{format=hang,justification=raggedright}|
\cap{\captionsetup{format=hang,justification=raggedright}}

|\captionsetup{format=hang,justification=Centering}|
\cap{\captionsetup{format=hang,justification=Centering}}

|\captionsetup{format=hang,justification=RaggedLeft}|
\cap{\captionsetup{format=hang,justification=RaggedLeft}}

|\captionsetup{format=hang,justification=RaggedRight}|
\cap{\captionsetup{format=hang,justification=RaggedRight}}

|\DeclareCaptionJustification{verbatim}{\@verbatim}|\\
|\captionsetup{justification=verbatim,singlelinecheck=0}|
\makeatletter
\DeclareCaptionJustification{verbatim}{\@verbatim}
\makeatother
\cap{\captionsetup{justification=verbatim,singlelinecheck=0}}

\captionsetup{singlelinecheck=1}


%************************************%
%*                                  *%
%* Fonts                            *%
%*                                  *%
%************************************%

\clearpage
\section{Fonts}

|\captionsetup{font=scriptsize}|
\cap{\captionsetup{font=scriptsize}}

|\captionsetup{font=footnotesize}|
\cap{\captionsetup{font=footnotesize}}

|\captionsetup{font=small}|
\cap{\captionsetup{font=small}}

|\captionsetup{font=normalsize}|
\cap{\captionsetup{font=normalsize}}

|\captionsetup{font=large}|
\cap{\captionsetup{font=large}}

|\captionsetup{font=Large}|
\cap{\captionsetup{font=Large}}

|\DeclareCaptionFont{LARGE}{\LARGE}|\\
|\captionsetup{size=LARGE}|
\DeclareCaptionFont{LARGE}{\LARGE}
\cap{\captionsetup{size=LARGE}}

|\captionsetup{textfont=up}|
\cap{\captionsetup{textfont=up}}

|\captionsetup{textfont=it}|
\cap{\captionsetup{textfont=it}}

|\captionsetup{textfont=sl}|
\cap{\captionsetup{textfont=sl}}

|\captionsetup{labelfont=sc}|
\cap{\captionsetup{labelfont=sc}}

|\captionsetup{labelfont=md}|
\cap{\captionsetup{labelfont=md}}

|\captionsetup{labelfont=bf}|
\cap{\captionsetup{labelfont=bf}}

|\captionsetup{font=rm}|
\cap{\captionsetup{font=rm}}

|\captionsetup{font=sf}|
\cap{\captionsetup{font=sf}}

|\captionsetup{font=tt}|
\cap{\captionsetup{font=tt}}

% Doppeltes Leerzeichen ist kein Bug, sondern wird
% nach ':' etc. automatisch so gemacht:
% {\ttfamily Figure 131: Short caption}

|\captionsetup{textfont=it,labelfont=bf,font=sf}|
\cap{\captionsetup{textfont=it,labelfont=bf,font=sf}}

|\captionsetup{labelfont={large,bf,sl}}|
\cap{\captionsetup{labelfont={large,bf,sl}}}


%************************************%
%*                                  *%
%* Margins                          *%
%*                                  *%
%************************************%

\clearpage
\section{Margins}

\hrule

|\captionsetup{margin=\leftmargini}|
\cap{\captionsetup{margin=\leftmargini}}

\hrule

|\captionsetup{margin={10pt,20pt}}|
\cap{\captionsetup{margin={10pt,20pt}}}

\hrule

|\captionsetup{margin={20pt,10pt}}|
\cap{\captionsetup{margin={20pt,10pt}}}

\hrule

|\captionsetup{width=4in}|
\cap{\captionsetup{width=4in}}

\hrule

|\captionsetup{margin={.5cm,2cm},width=3in}|
\cap{\captionsetup{margin={.5cm,2cm},width=3in}}

\hrule

%************************************%
%*                                  *%
%* Local settings                   *%
%*                                  *%
%************************************%

\clearpage
\section{Local settings}

%|\captionsetup[figure]{style=italic,name=FIGURE}|\\
|\captionsetup[figure]{style=italic}|\\
|\captionsetup[table]{position=top}|

%\captionsetup[figure]{style=italic,name=FIGURE}
\captionsetup[figure]{style=italic}
\captionsetup[table]{position=top}

\begin{figure}[!ht]
\centering Text
\caption{This is a figure}
\end{figure}

\begin{table}[!ht]
\caption{This is a table}
\centering\begin{tabular}{ll}
A & B \\
C & D \\
\end{tabular}
\end{table}

|\clearcaptionsetup{figure}|\\
|\clearcaptionsetup{table}|

\clearcaptionsetup{figure}
\clearcaptionsetup{table}

\begin{figure}[!ht]
\centering Text
\caption{This is a figure}
\end{figure}

\begin{table}[!ht]
\caption{This is a table}
\centering\begin{tabular}{ll}
A & B \\
C & D \\
\end{tabular}
\end{table}


%************************************%
%*                                  *%
%* Hooks                            *%
%*                                  *%
%************************************%

\clearpage
\section{Hooks}

|\AtBeginCaption{$<$}|
|\AtEndCaption{$>$}|

\cap{\AtBeginCaption{$<$}\AtEndCaption{$>$}}


%************************************%
%*                                  *%
%* Multiparagraph captions          *%
%*                                  *%
%************************************%

\clearpage
\section{Multiparagraph captions}
\begingroup

\captionsetup{justification=centerfirst}
\captionof{figure}[]{A caption that runs over more than one line
    to show the effect of the centerfirst keyword.

   It even has a second paragraph, which is kind of
   problematical with this type of layout.}

\captionsetup{justification=centerlast}
\captionof{figure}[]{A caption that runs over more than one line
    to show the effect of the centerlast keyword.

   It even has a second paragraph, which is kind of
   problematical with this type of layout.}

\captionsetup{justification=justified}

\newcommand*\mptest[1]{%
\noindent{\ttfamily#1}\par
\captionsetup{#1}
\captionof{figure}[]{Short caption}
\captionof{figure}[]{%
  This is some text. This is some text. This is some text.
  This is some text. This is some text. This is some text.
  This is some text. This is some text. This is some text.\\
  This is some text after a hard line break.}
\captionof{figure}[]{%
  Long long long long long long long long
  long long long long long long long long long long long long
  long long long long long long long long long long long caption

  Long long long long long long long long%\\
  long long long long long long long long long long long long
  long long long long long long long long long long long caption

  Long long long long long long long long
  long long long long long long long long long long long long
  long long long long long long long long long long long caption}
}

\mptest{format=normal,indention=0cm}

\mptest{format=normal,indention=.5cm}

\mptest{format=normal,indention=-.5cm}

\mptest{format=hang,indention=0cm}

\mptest{format=hang,indention=.5cm}

\mptest{format=hang,indention=-.5cm}

\mptest{format=normal,indention=0cm,hangindent=.5cm}
% = format=default-from-steven,indention=.5cm

\mptest{format=hang,indention=0cm,hangindent=-.5cm}
% = format=hang-from-steven,indention=-.5cm

\mptest{format=hang,indention=0cm,hangindent=.5cm}
% = format=hang-from-steven,indention=.5cm

\endgroup


%********************************************%
%*                                          *%
%* The \caption and \ContinuedFloat command *%
%*                                          *%
%********************************************%

\clearpage
\section{The \cs{caption} and \cs{ContinuedFloat} command}
\captionsetup{position=auto,labelseparator=visiblespace}

\begin{table}[!ht]
\caption{This is a table typeset with \cs{caption}}
\centering\begin{tabular}{ll}
A & B \\
C & D \\
\end{tabular}
\end{table}

\begin{table}[!ht]
\caption*{This is a table typeset with \cs{caption*}}
\centering\begin{tabular}{ll}
A & B \\
C & D \\
\end{tabular}
\end{table}

\begin{table}[!ht]
\caption[]{This is a table typeset with \cs{caption[]}}
\centering\begin{tabular}{ll}
A & B \\
C & D \\
\end{tabular}
\end{table}

\iffalse
\begin{table}[!ht]
\caption*[Huba!]{This is a table typeset with \cs{caption*[Huba!]}}
\centering\begin{tabular}{ll}
A & B \\
C & D \\
\end{tabular}
\end{table}
\fi

\begin{table}[!ht]
\caption{} % BUG! PRODUCES NO LOT!
\centering\begin{tabular}{ll}
A & B \\
C & D \\
\end{tabular}
\end{table}

\begin{table}[!ht]
\caption*{}
\centering\begin{tabular}{ll}
A & B \\
C & D \\
\end{tabular}
\end{table}

\begin{table}[!ht]
\caption{This is a table}
\centering\begin{tabular}{ll}
A & B \\
C & D \\
\end{tabular}
\end{table}

\begin{table}[!ht]
\ContinuedFloat
\caption{This is a table (cont.)}
\centering\begin{tabular}{ll}
A & B \\
C & D \\
\end{tabular}
\end{table}

\begingroup
  \let\section\oldsection
  \listoftables
\endgroup

\captionsetup{style=default,position=default,aboveskip=10pt,belowskip=0pt}


%************************************%
%*                                  *%
%* Test of minipage flag            *%
%*                                  *%
%************************************%

\clearpage
\section{Test of minipage flag}

\begin{minipage}\linewidth
\makeatletter
\if@minipage
\typeout{@minipage is set}
\else
\error{minipage}
\fi
\makeatother
\captionof{figure}{A minipage}
\makeatletter
\if@minipage
\error{minipage}
\else
\typeout{@minipage is not set}
\fi
\makeatother

A minipage
\end{minipage}


%************************************%
%*                                  *%
%* Packages                         *%
%*                                  *%
%************************************%

\captionsetup{font=sl,labelfont=bf,labelsep=period}
\DeleteShortVerb{\|}


%************************************%
%*                                  *%
%* The subfig package               *%
%*                                  *%
%************************************%

\clearpage
\section{The subfig package}

\ifundefined{subfloat}
---\par
(Please use \verb|testsub.tex| to test the subfig package.)
\else

\begingroup

\newcommand\LCap{A longer caption with more text}
\newcommand\FIG{\fbox{\parbox{.4\textwidth}{\strut}}}

\begin{figure}[ht]  \centering
  \subfloat[Short caption]{\FIG}
  \subfloat[\LCap]{\FIG}
  \caption{Default subfigures}
\end{figure}

\captionsetup[subfloat]{format=hang,textfont=it,labelfont={rm,bf},labelformat=simple,labelsep=period,margin=5pt,justification=raggedright}

\begin{figure}[ht]  \centering
  \subfloat[Short caption]{\FIG}
  \subfloat[\LCap]{\FIG}
  \caption{Customized subfigures}
\end{figure}

\endgroup

\fi


%************************************%
%*                                  *%
%* The float package                *%
%*                                  *%
%************************************%

\clearpage
\section{The float package}

\makeatletter

\newcommand\thecaptionposition{%
  \caption@iftop{top}{bottom}}

\newcommand\thefloatstyle{%
  \ifx\caption@fixfloat@c\undefined
    \caption@fst
  \else
    \begingroup
      \let\caption@fixfloat@c\@gobble
      \let\caption@setstyle\@gobble
      \let\caption@setfloatposition\relax
      \let\caption@settype\@gobble
      \caption@setfloattype{\@captype}%
      \caption@fst
    \endgroup
  \fi}

\makeatother

\ifundefined{floatstyle}
---
\else
\begingroup

\ifundefined{floatrow}
  \newcommand\floatinfo{%
    (float style: `\thefloatstyle', caption position: \thecaptionposition)}
\else
  \newcommand\floatinfo{%
    (caption position: \thecaptionposition)}
\fi

\begin{Program}[!ht]
\begin{verbatim}
#include <stdio.h>

int main(int argc, char **argv)
{
       int i;
       for (i = 0; i < argc; ++i)
               printf("argv[%d] = %s\n", i, argv[i]);
       return 0;
}
\end{verbatim}
\caption{The first program. This hasn't got anything to do with the package
   but is included as an example. Note the \texttt{ruled} float style.%
   \label{prog1.1}}
%\floatinfo
\end{Program}

\captionsetup[ruled]{margin=10pt} % new feature v3.0f

\begin{Program}[!ht]
\begin{verbatim}
#include <stdio.h>

int main(int argc, char **argv)
{
       int i;
       for (i = 0; i < argc; ++i)
               printf("argv[%d] = %s\n", i, argv[i]);
       return 0;
}
\end{verbatim}
\caption{The first program. This hasn't got anything to do with the package
   but is included as an example. Note the extra margin of \texttt{10pt}.}
%\floatinfo
\end{Program}

\clearcaptionsetup{ruled}

\begin{Program}[!ht]
\begin{verbatim}
#include <stdio.h>

int main(int argc, char **argv)
{
       int i;
       for (i = 0; i < argc; ++i)
               printf("argv[%d] = %s\n", i, argv[i]);
       return 0;
}
\end{verbatim}
\captionsetup{textfont=sf} % new feature v3.0f
\caption{The first program. This hasn't got anything to do with the package
   but is included as an example. Note the serif text font.}
%\floatinfo
\end{Program}

%\clearpage

\MakeShortVerb{\|}

\begin{Example}[H]
\begin{verse}
|\floatstyle{ruled}|\\
|\newfloat{Program}{tbp}{lop}[section]|\\
\dots\ loads o' stuff \dots\\
|\begin{Program}|\\
|\begin{verbatim}|\\
\dots\ program text \dots\\
|\end{verbatim}|\\
|\caption{|\dots\ caption \dots|}|\\
|\end{Program}|
\end{verse}
\caption{This is another silly floating Example. Except that this one
  doesn't actually float because it uses the {\tt[H]} optional parameter
  to appear \textbf{Here}. (Gotcha.)}
%\floatinfo
\end{Example}

\DeleteShortVerb{\|}

%\clearpage

\floatstyle{boxed}
\restylefloat{table}

\begin{table}[!ht] \def\B#1{$\displaystyle{n\choose#1}$}
\begin{center} \begin{tabular}{c|cccccccc}
$n$&\B0&\B1&\B2&\B3&\B4&\B5&\B6&\B7\\ \hline
 0 & 1\\
 1 & 1&1\\
 2 & 1&2&1\\
 3 & 1&3&3&1\\
 4 & 1&4&6&4&1\\
 5 & 1&5&10&10&5&1\\
 6 & 1&6&15&20&15&6&1\\
 7 & 1&7&21&35&35&21&7&1
\end{tabular} \end{center}
\caption{Pascal's triangle. This is a re-styled \LaTeX\ \texttt{table}.%
  \label{table1}}
%\floatinfo
\end{table}

\captionsetup[boxed]{skip=10pt}

\begin{table}[!ht] \def\B#1{$\displaystyle{n\choose#1}$}
\begin{center} \begin{tabular}{c|cccccccc}
$n$&\B0&\B1&\B2&\B3&\B4&\B5&\B6&\B7\\ \hline
 0 & 1\\
 1 & 1&1\\
 2 & 1&2&1\\
 3 & 1&3&3&1\\
 4 & 1&4&6&4&1\\
 5 & 1&5&10&10&5&1\\
 6 & 1&6&15&20&15&6&1\\
 7 & 1&7&21&35&35&21&7&1
\end{tabular} \end{center}
\caption{Pascal's triangle. This is a re-styled \LaTeX\ \texttt{table} with \texttt{skip=10pt}.}
%\floatinfo
\end{table}

\clearpage

\floatstyle{plaintop}
\restylefloat{table}

\begin{table}[!ht]
\begin{center}
\hrule
Plaintop float
\hrule
\end{center}
\caption{This is a re-styled \LaTeX\ \texttt{table}.%
  \label{table2}}
%\floatinfo
\end{table}

%\clearpage

\begingroup
\hrule
\captionof{Program}{Caption of Programm}
\hrule
\endgroup

\begingroup
\hrule
\captionof{Example}{Caption of Example}
\hrule
\endgroup

\begingroup
\hrule
\captionof{figure}{Caption of figure}
\hrule
\endgroup

\begingroup
\hrule
\captionof{table}{Caption of table}
\hrule
\endgroup

\endgroup

\clearpage
Same with \verb|\restylefloat*|\ \ldots

\begingroup

\floatstyle{ruled}
\restylefloat*{Program}
\floatstyle{plain}
\restylefloat*{Example}

\begin{Program}[!ht]
\begin{verbatim}
#include <stdio.h>

int main(int argc, char **argv)
{
       int i;
       for (i = 0; i < argc; ++i)
               printf("argv[%d] = %s\n", i, argv[i]);
       return 0;
}
\end{verbatim}
\caption{The first program. This hasn't got anything to do with the package
   but is included as an example. Note the \texttt{ruled} float style.}
(The position of the caption is \thecaptionposition)
\end{Program}

%\clearpage

\MakeShortVerb{\|}

\begin{Example}[H]
\begin{verse}
|\floatstyle{ruled}|\\
|\newfloat{Program}{tbp}{lop}[section]|\\
\dots\ loads o' stuff \dots\\
|\begin{Program}|\\
|\begin{verbatim}|\\
\dots\ program text \dots\\
|\end{verbatim}|\\
|\caption{|\dots\ caption \dots|}|\\
|\end{Program}|
\end{verse}
%\StartTracing
\caption{This is another silly floating Example. Except that this one
  doesn't actually float because it uses the {\tt[H]} optional parameter
  to appear \textbf{Here}. (Gotcha.)}
%\StopTracing
(The position of the caption is \thecaptionposition)
\end{Example}

\DeleteShortVerb{\|}

%\clearpage

\floatstyle{boxed}
\restylefloat*{table}

\begin{table}[!ht] \def\B#1{$\displaystyle{n\choose#1}$}
\begin{center} \begin{tabular}{c|cccccccc}
$n$&\B0&\B1&\B2&\B3&\B4&\B5&\B6&\B7\\ \hline
 0 & 1\\
 1 & 1&1\\
 2 & 1&2&1\\
 3 & 1&3&3&1\\
 4 & 1&4&6&4&1\\
 5 & 1&5&10&10&5&1\\
 6 & 1&6&15&20&15&6&1\\
 7 & 1&7&21&35&35&21&7&1
\end{tabular} \end{center}
\caption{Pascal's triangle. This is a re-styled \LaTeX\ \texttt{table}.}
(The position of the caption is \thecaptionposition)
\end{table}

\iffalse
\floatstyle{plaintop}
\restylefloat*{table}

\begin{table}[!ht]
\begin{center}
Plaintop float
\end{center}
\caption{This is a re-styled \LaTeX\ \texttt{table}.%
  \label{table4}}
(The position of the caption is \thecaptionposition)
\end{table}
\fi

\clearpage

\begingroup
\hrule
\captionof{Program}{Caption of Programm}
\hrule
\endgroup

\begingroup
\hrule
\captionof{Example}{Caption of Example}
\hrule
\endgroup

\begingroup
\hrule
\captionof{figure}{Caption of figure}
\hrule
\endgroup

\begingroup
\hrule
\captionof{table}{Caption of table}
\hrule
\endgroup

\endgroup

\verb|\caption*{Caption}|
\begingroup
%\ifundefined{subfloat}
  \captionof*{Example}{Caption}
%\else\par
%  ******* BUG: \verb|\captionof*| does not work with \textsf{subfig}!
%  (\verb|\ContinuedFloat| does not work inside environments defined using
%   the \textsf{float} package) *******
%\fi
\endgroup
\begingroup
\captionof*{figure}{Caption}
\endgroup

\verb|\caption{}|
\begingroup
\captionof{Example}{}
\endgroup
\begingroup
\captionof{figure}{}
\endgroup

\fi


%************************************%
%*                                  *%
%* The listings package             *%
%*                                  *%
%************************************%

\clearpage
\section{The listings package}

\ifundefined{lstlisting}
---
\else
\begingroup

   \begin{lstlisting}[caption={Useless code},label=useless]
   for i:=maxint to 0 do
   begin
       { do nothing }
   end;
   \end{lstlisting}
   \begin{lstlisting}[title={`Caption' without label}]
   for i:=maxint to 0 do
   begin
       { do nothing }
   end;
   \end{lstlisting}

%   \begin{lstlisting}[caption={ },label=useless]
%   for i:=maxint to 0 do
%   begin
%       { do nothing }
%   end;
%   \end{lstlisting}

\MakeShortVerb{\|}
|\captionsetup{labelformat=parens,labelsep=space}|\\
|\captionsetup[lstlisting]{textfont={up,md,rm}}|
\captionsetup{labelformat=parens,labelsep=space}
\captionsetup[lstlisting]{textfont={up,md,rm}}
\DeleteShortVerb{\|}

   \begin{lstlisting}[caption={Useless code},label=useless2]
   for i:=maxint to 0 do
   begin
       { do nothing }
   end;
   \end{lstlisting}
   \begin{lstlisting}[title={`Caption' without label}]
   for i:=maxint to 0 do
   begin
       { do nothing }
   end;
   \end{lstlisting}

%   \begin{lstlisting}[caption={},label=useless]
%   for i:=maxint to 0 do
%   begin
%       { do nothing }
%   end;
%   \end{lstlisting}

\endgroup
\fi


%************************************%
%*                                  *%
%* The longtable package            *%
%*                                  *%
%************************************%

\clearpage
\section{The longtable package}

\ifundefined{longtable}
---
\else
\begingroup

\makeatletter
\let\LT@makecaption\Orig@LT@makecaption
\makeatother

\begin{longtable}{lll}
\caption{The ISOGRK3 entity set (longtable code)}\\
\bfseries Entity&\bfseries  Unicode Name&\bfseries  Unicode\\
\hline
\endfirsthead
\bfseries Entity&\bfseries  Unicode Name&\bfseries  Unicode\\
\hline
\endhead
\hline
\multicolumn{3}{r}{\emph{Continued on next page}}
\endfoot
\hline
\endlastfoot
alpha              & GREEK SMALL LETTER ALPHA            & 03B1\\
beta               & GREEK SMALL LETTER BETA             & 03B2\\
chi                & GREEK SMALL LETTER CHI              & 03C7\\
\end{longtable}

\endgroup\begingroup

\captionsetup{position=default}

\begin{longtable}{lll}
\caption{The ISOGRK3 entity set (caption code)}\\
\bfseries Entity&\bfseries  Unicode Name&\bfseries  Unicode\\
\hline
\endfirsthead
\bfseries Entity&\bfseries  Unicode Name&\bfseries  Unicode\\
\hline
\endhead
\hline
\multicolumn{3}{r}{\emph{Continued on next page}}
\endfoot
\hline
\endlastfoot
alpha              & GREEK SMALL LETTER ALPHA            & 03B1\\
beta               & GREEK SMALL LETTER BETA             & 03B2\\
chi                & GREEK SMALL LETTER CHI              & 03C7\\
\end{longtable}

\restylefloat{table} % This should not affect the longtable!
\captionsetup[longtable]{labelsep=colon}

\begin{longtable}{ll}
A & B \\
\caption{\longtext}\\
C & D \\
\end{longtable}

\captionsetup[table]{margin=0pt} % should be overwritten by \LTcapwidth

Setting \texttt{LTcapwidth=3in}\ldots
\LTcapwidth=3in

\begin{longtable}{ll}
\caption{\longtext}\\
A & B \\
C & D \\
\end{longtable}

Setting \texttt{LTcapwidth=3.999in}\ldots
\LTcapwidth=3.999in

\begin{longtable}{ll}
\caption{\longtext}\\
A & B \\
C & D \\
\end{longtable}

Setting \texttt{LTcapwidth=3.999in} and \verb|\captionsetup[longtable]{width=3in}|\ldots
\LTcapwidth=3.999in
\captionsetup[longtable]{width=3in}

\begin{longtable}{ll}
\caption{\longtext}\\
A & B \\
C & D \\
\end{longtable}

Setting \texttt{LTcapwidth=4in} again\ldots
\LTcapwidth=4in
\clearcaptionsetup{longtable}

\begin{longtable}{ll}
\caption{\longtext}\\
A & B \\
C & D \\
\end{longtable}

\clearpage

Caption position is now \texttt{default}:
\captionsetup{position=default,debug=0}

\begin{longtable}{ll}
A & B \\
\caption{\shorttext}\\
C & D \\
\end{longtable}

Caption position is now \texttt{top}:
\captionsetup{position=top}

\begin{longtable}{ll}
A & B \\
\caption{\shorttext}\\
C & D \\
\end{longtable}

Caption position is now \texttt{bottom}:
\captionsetup{position=bottom}

\begin{longtable}{ll}
A & B \\
\caption{\shorttext}\\
C & D \\
\end{longtable}

Caption position is now \texttt{auto}:
\captionsetup{position=auto}

\begin{longtable}{ll}
A & B \\
\caption{\shorttext}\\
C & D \\
\end{longtable}

\iffalse
\clearpage

Longtable column position is now \texttt{l}:
\def\longtablecolumnpos{l}

\begin{longtable}{ll}
A & B \\
\caption{\shorttext}\\
C & D \\
\end{longtable}

\begin{longtable}[l]{ll}
A & B \\
\caption{\shorttext}\\
C & D \\
\end{longtable}

\begingroup
\def\makelongtablecaption#1{\hbox to 0pt{\parbox[t]\linewidth{#1}\hss}}%
\begin{longtable}[l]{ll}
A & B \\
\caption{\shorttext}\\
C & D \\
\end{longtable}
\captionsetup{singlelinecheck=0}
\begin{longtable}[l]{ll}
A & B \\
\caption{\shorttext}\\
C & D \\
\end{longtable}
\endgroup

Longtable column position is now \texttt{c}:
\def\longtablecolumnpos{c}

\begin{longtable}{ll}
A & B \\
\caption{\shorttext}\\
C & D \\
\end{longtable}

Longtable column position is now \texttt{r}:
\def\longtablecolumnpos{r}

\begin{longtable}{ll}
A & B \\
\caption{\shorttext}\\
C & D \\
\end{longtable}

\def\longtablecolumnpos{c}

\fi

\verb|\caption*{Caption}|
\begin{longtable}{ll}
\caption*{\shorttext}\\
A & B \\
C & D \\
\end{longtable}

\verb|\caption{}|
\begin{longtable}{ll}
\caption{}\\
A & B \\
C & D \\
\end{longtable}

\endgroup
\fi


%************************************%
%*                                  *%
%* The rotating package             *%
%*                                  *%
%************************************%

\clearpage
\section{The rotating package}

\ifundefined{sidewaystable}
---
\else

\begin{sidewaystable}
\centering
\caption{This is a narrow  table, which should be centred vertically
on the final page.\label{rotfloat1}}
  \begin{tabular}{|ll|}
\hline
    a & b \\
    c & d \\
    e & f \\
    g & h \\
    i & j \\
\hline
  \end{tabular}
\end{sidewaystable}

\begin{sidewaystable}
\centering
\begin{tabular}{|llllllllp{1in}lp{1in}|}
\hline
Context   &Length   &Breadth/   &Depth   &Profile   &Pottery   &Flint   &Animal   &Stone   &Other    &C14 Dates \\
  &         &Diameter   &        &          &          &        & 
Bones&&&\\
\hline
&&&&&&&&&&\\
\multicolumn{10}{|l}{\bf Grooved Ware}&\\
784       &---        &0.9m       &0.18m   &Sloping U &P1       &$\times$46  &  $\times$8      &&       $\times$2 bone&  2150$\pm$ 100 BC\\
785       &---        &1.00m      &0.12    &Sloping U &P2--4    &$\times$23  &  $\times$21     & Hammerstone &---&---\\
962       &---        &1.37m      &0.20m   &Sloping U &P5--6    &$\times$48  &  $\times$57*    & ---&     ---&1990 $\pm$ 80 BC (Layer 4) 1870 $\pm$90 BC (Layer 1)\\
983       &0.83m      &0.73m      &0.25m   &Stepped U &---      &$\times$18  &  $\times$8      & ---& Fired clay&---\\
&&&&&&&&&&\\
\multicolumn{10}{|l}{\bf Beaker}&\\
552       &---        &0.68m      &0.12m   &Saucer    &P7--14   &---           & ---       & ---       &---        &---\\
790       &---        &0.60m      &0.25m   &U         &P15      &$\times$12    & ---       & Quartzite-lump&---    &---\\
794       &2.89m      &0.75m      &0.25m   &Irreg.    &P16      &$\times$3     & ---       & ---       &---        &---\\
\hline
\end{tabular}
 
\caption[Grooved Ware and Beaker Features, their Finds and
Radiocarbon Dates]{Grooved Ware and Beaker Features, their
Finds and Radiocarbon Dates; For a breakdown of the Pottery
Assemblages see Tables I and III; for
the Flints see Tables II and IV; for the
Animal Bones see Table V.}\label{rotfloat2}
\end{sidewaystable}

\begin{table}
\centering
\hbox{
\rotcaption{Minimum number of individuals; effect of rotating table
and caption separately}\label{rotfloat3}%
\begin{sideways}
\begin{tabular}[t]{cccccccccp{1cm}}
\hline
Phase&Total&Cattle&Sheep&Pig&Red Deer&Horse&Dog&Goat&Other\\
\hline
&1121&54&12&32&1&1&1&1&1 polecat\\
3&8255&58&6&35&1&1&1&1&1 roe deer, 1 hare, 1 cat, 1 otter\\
4&543&45&6&45&4&1&1&---&---\\
\hline
&9919&157&24&112&6&3&3&2&5\\
\hline
\end{tabular}
\end{sideways}
}
\end{table}

\fi


%************************************%
%*                                  *%
%* The sidecap package              *%
%*                                  *%
%************************************%

\clearpage
\section{The sidecap package}

\ifundefined{SCfigure}
---
\else
\begingroup

\captionsetup[SCtable]{labelsep=colon}

\newcommand{\FIG}[2][]{%
  \begingroup
    \def\xxx{#1}%
    \ifx\xxx\empty
      \setlength{\unitlength}{\linewidth}%
      \addtolength{\unitlength}{-2\fboxrule}%
      \setlength{\unitlength}{.1\unitlength}%
    \else  
      \setlength{\unitlength}{#1}%
    \fi% 
    \setlength{\fboxsep}{0pt}%
    \ifcase#2\relax
    \or%
      \fbox{%
	\begin{picture}(4,6)%
	  \put(1,5){\circle{1}}%
	  \put(3,5){\circle{1}}%
	  \put(2,3){\circle{1}}%
	  \put(1,1){\circle{1}}%
	  \put(3,1){\circle{1}}%     
	\end{picture}}%
    \or%	
      \fbox{%
	\begin{picture}(10,4)%
	  \put(1,3){\circle{1}}%
	  \put(9,3){\circle{1}}%
	  \put(5,2){\circle{1}}%
	  \put(1,1){\circle{1}}%
	  \put(9,1){\circle{1}}%     
	\end{picture}}%
    \or% (doesn't make a difference to use [b] or [t] or [c] here !!!)
      \begin{minipage}{5\unitlength}%
        text text text text text text text text
	text text text text text text text text
	text text text text text text text text
	text text text text text text text text
	text text text text text text text text
	text text text text text text text text
	text text text text text text text text
	text text text text text text text text.
      \end{minipage}%	
    \or% (doesn't make a difference to use [b] or [t] or [c] here !!!)
      \begin{minipage}{5\unitlength}%
        \strut
        text text text text text text text text
	text text text text text text text text
	text text text text text text text text
	text text text text text text text text
	text text text text text text text text
	text text text text text text text text
	text text text text text text text text
	text text text text text text text text.
	\unskip\strut
      \end{minipage}%	
    \fi	
  \endgroup  
}

\newcommand{\TABi}{%
  \begin{tabular}{@{} | c | c c c c c | @{}} \hline
    \multicolumn{1}{| c} {1} 
       &  2 &  3 &  4 &  5 &  6 \\ \cline{2-6}
     2 &  4 &  6 &  8 & 10 & 12 \\
     3 &  6 &  9 & 12 & 15 & 18 \\     
     4 &  8 & 12 & 16 & 20 & 24 \\
     5 & 10 & 15 & 20 & 25 & 30 \\
     6 & 12 & 18 & 24 & 30 & 36 \\
     7 & 14 & 21 & 28 & 35 & 42 \\
     8 & 16 & 24 & 32 & 40 & 48 \\
     9 & 18 & 27 & 36 & 45 & 54 \\
    10 & 20 & 30 & 40 & 50 & 60 \\
    11 & 22 & 33 & 44 & 55 & 66 \\
    12 & 24 & 36 & 48 & 60 & 72 \\
    13 & 26 & 39 & 52 & 65 & 78 \\
    14 & 28 & 42 & 56 & 70 & 84 \\
    \hline
  \end{tabular}%
}

\newcommand{\CAPi}{%
  The dazed brown fox quickly gave 12345--67890 jumps!
  The dazed brown fox quickly gave 12345--67890 jumps!
}
  
\newcommand{\CAPii}{%
  The dazed brown fox quickly gave 12345--67890 jumps!
  The dazed brown fox quickly gave 12345--67890 jumps!
  The dazed brown fox quickly gave 12345--67890 jumps!
  The dazed brown fox quickly gave 12345--67890 jumps!
}
  
\newcommand{\SHORTCAPi}{%
  The dazed brown fox etc.}  
  
\newcommand*{\MARKER}
  {\noindent\strut\vrule\hrulefill~text area 
   (\ifthenelse{\isodd{\thepage}}{odd}{even} page)~\hrulefill\vrule%
    \marginpar{\strut\vrule\hrulefill~margin area~\hrulefill\vrule}}  

%\section*{The \texttt{sidecap} package}

First, a standard \LaTeX\ \texttt{figure} environment
(fig.~\ref{fig:standard}).
Since width of the body of the figure is considerably less than
the document's \verb|\textwidth|, there is much space wasted.

\begin{figure}[h]
  \centering
  \FIG[1cm]{1}%
  \caption[\SHORTCAPi]{\CAPi}
  \label{fig:standard}
\end{figure}

For such cases, the \textsf{sidecap} package provides the
\verb|SCfigure| and \verb|SCtable| environments to put the
caption beside the body of the floating environment.

In this document, we used the defaults of the \textsf{sidecap} package:
%
\begin{verbatim}
  \usepackage{sidecap}
\end{verbatim}
%
In the text, we used the optional arguments for the following figures
(fig.~\ref{fig:sc1}, \ref{fig:sc2}, and \ref{fig:sc3}).
%
\begin{verbatim}
  \begin{SCfigure} ...
  \begin{SCfigure}[1.2][bhp] ...
  \begin{SCfigure}[][bhp] ...
\end{verbatim}
%
In the first case, no optional arguments were given, so the defaults
(\verb|[1.0][tbp]|) are used. I.e., the caption will be as wide as
the body of the float; and the `float placement specifier' is the same
as for \LaTeX's standard floats (fig.~\ref{fig:sc1}).
In the second case, both optional arguments are given to change the
caption width as well as the placement options (fig.~\ref{fig:sc2}).
To change only the placement, just leave the first optional argument
empty (fig.~\ref{fig:sc3}).

\begin{SCfigure}
  \FIG[1cm]{1}%
%%%\captionsetup{font=large}
  \caption[\SHORTCAPi]{\CAPi}
  \label{fig:sc1}
\end{SCfigure}

\begin{SCfigure}[1.2][bhp]
  \FIG[1cm]{1}%
  \caption[\SHORTCAPi]{\CAPi}
  \label{fig:sc2}
\end{SCfigure}

\begin{SCfigure}[][bhp]
  \FIG[1cm]{1}%
  \caption[\SHORTCAPi]{\CAPi}
  \label{fig:sc3}
\end{SCfigure}

\clearpage

Now, a standard \LaTeX\ \texttt{table} environment
(tab.~\ref{tab:standard}).

\begin{table}[h]
  \centering
  \caption[\SHORTCAPi]{\CAPi}
  \TABi%
  \label{tab:standard}
\end{table}

The only difference between the \verb|SCfigure| and the
\verb|SCtable| environments is the vertical position of the
caption with respect to the body of the float:
The former places the caption bottom-aligned, the latter top-aligned
(compare fig.~\ref{fig:sc1} to tab.~\ref{tab:sc1}).

Also, like standard \LaTeX, the \textsf{sidecap} package doesn't care
about the contents of the float (cf.\ tabs.~\ref{tab:standard} and
\ref{tab:sc1}).

\begin{SCtable}
  \caption[\SHORTCAPi]{\CAPi}  
  \FIG[1cm]{1}%
  \label{tab:sc1}
\end{SCtable}

If you want to use all the remaining space for the caption, just
specify a sufficiently large factor in the first optional argument.
E.g., for tab.~\ref{tab:sc2} we used the following:
%
\begin{verbatim}
  \begin{SCtable}[50] ...
\end{verbatim}
%
Note, however, that \TeX\ may complain
%
\begin{verbatim}
  ! Dimension too large.
\end{verbatim}
%
if the factor times the width of the body is huge enough to exceed
the range \TeX\ is able to represent in a length register.
%
(Max.\ \verb|\maxdimen|%
${}={}$\the\maxdimen
${}\approx 5.76$\,{meters}).

% VOODOO:
%\setlength
%\mylength=\maxdimen
%\mylength=-\mylength
%\the\mylength

\begin{SCtable}[50]
%  \caption[\SHORTCAPi]{\CAPi}
\captionsetup{font=,labelfont=sf}
\captionsetup{margin=10pt}
  \caption[]{\CAPi}
  \label{tab:sc2}
%  \TABi%
  \FIG[1cm]{1}%
\end{SCtable}

\verb|\caption*{Caption}|
\begin{SCfigure}[][bhp]
  \FIG[1cm]{1}%
  \caption*{\CAPi}
\end{SCfigure}

\verb|\caption{}|
\begin{SCfigure}[][bhp]
  \FIG[1cm]{1}%
  \caption{}
\end{SCfigure}

\captionsetup[figure]{margin=5cm}

\begin{SCfigure}[50][bhp]%
  \fbox{\parbox{.4\textwidth}{%
  Setting a figure margin
  should have no effect
  on sidecaps\ldots}}%
  \caption{\CAPi}
\end{SCfigure}

\captionsetup[SCfigure]{margin=2cm}

\begin{SCfigure}[50][bhp]%
  \fbox{\parbox{.4\textwidth}{%
  \ldots unless we set a SCfigure
  margin!}}%
  \caption{\CAPi}
\end{SCfigure}

\endgroup
\fi


%************************************%
%*                                  *%
%* The supertabular package         *%
%*                                  *%
%************************************%

\clearpage
\section{The supertabular package}

\ifundefined{supertabular}
---
\else

\setlength\parskip{0pt}  % \tablecaption reagiert allergisch auf andere \parskips...
\captionsetup{debug=2}

\begingroup
\makeatletter
\let\ST@caption\Orig@ST@caption
\let\@makecaption\Orig@makecaption
\makeatother

\begin{center}
\tablecaption{The ISOGRK3 entity set (supertabular code)}
\begin{supertabular}{lll}
\bfseries Entity&\bfseries  Unicode Name&\bfseries  Unicode\\
\hline
alpha              & GREEK SMALL LETTER ALPHA            & 03B1\\
beta               & GREEK SMALL LETTER BETA             & 03B2\\
chi                & GREEK SMALL LETTER CHI              & 03C7\\
\end{supertabular}
\end{center}

\begin{center}
\bottomcaption{The ISOGRK3 entity set (supertabular code)}
\begin{supertabular}{lll}
\bfseries Entity&\bfseries  Unicode Name&\bfseries  Unicode\\
\hline
alpha              & GREEK SMALL LETTER ALPHA            & 03B1\\
beta               & GREEK SMALL LETTER BETA             & 03B2\\
chi                & GREEK SMALL LETTER CHI              & 03C7\\
\end{supertabular}
\end{center}
\endgroup

\begingroup
\begin{center}
\tablecaption{The ISOGRK3 entity set (caption code)}
\begin{supertabular}{lll}
\bfseries Entity&\bfseries  Unicode Name&\bfseries  Unicode\\
\hline
alpha              & GREEK SMALL LETTER ALPHA            & 03B1\\
beta               & GREEK SMALL LETTER BETA             & 03B2\\
chi                & GREEK SMALL LETTER CHI              & 03C7\\
\end{supertabular}
\end{center}

\captionsetup[supertabular]{labelsep=colon}

\begin{center}
\bottomcaption{The ISOGRK3 entity set (caption code)}
\begin{supertabular}{lll}
\bfseries Entity&\bfseries  Unicode Name&\bfseries  Unicode\\
\hline
alpha              & GREEK SMALL LETTER ALPHA            & 03B1\\
beta               & GREEK SMALL LETTER BETA             & 03B2\\
chi                & GREEK SMALL LETTER CHI              & 03C7\\
\end{supertabular}
\end{center}
\endgroup

\captionsetup{debug=1}

\verb|\caption*{Caption}|
\begin{center}
\tablecaption*{The ISOGRK3 entity set}
\begin{supertabular}{lll}
\bfseries Entity&\bfseries  Unicode Name&\bfseries  Unicode\\
\hline
alpha              & GREEK SMALL LETTER ALPHA            & 03B1\\
beta               & GREEK SMALL LETTER BETA             & 03B2\\
chi                & GREEK SMALL LETTER CHI              & 03C7\\
\end{supertabular}
\end{center}

\verb|\caption{}|
\begin{center}
\tablecaption{}
\begin{supertabular}{lll}
\bfseries Entity&\bfseries  Unicode Name&\bfseries  Unicode\\
\hline
alpha              & GREEK SMALL LETTER ALPHA            & 03B1\\
beta               & GREEK SMALL LETTER BETA             & 03B2\\
chi                & GREEK SMALL LETTER CHI              & 03C7\\
\end{supertabular}
\end{center}

\fi


%************************************%
%*                                  *%
%* End of Packages                  *%
%*                                  *%
%************************************%

\MakeShortVerb{\|}
\captionsetup{style=default,position=default,aboveskip=10pt,belowskip=0pt}


%************************************%
%*                                  *%
%* Reported bugs                    *%
%*                                  *%
%************************************%

\clearpage
\section{Reported bugs}

%-----------------------------------------------------------------------------%

\captionsetup{format=hang}

\begin{figure}[!htb]
This figure is labeled correctly.
\caption{``A sobering thought, Eileen: What if, right at this very moment I
*am* living up to my full potential?'' -- Jane Wagner\label{xxx}}
\end{figure}

% =>
% \@writefile{toc}{\contentsline {section}{\numberline {21}Reported bugs}{39}}
% \newlabel{xxx}{{174}{39}}
% \@writefile{lof}{\contentsline {figure}{\numberline {174}{\ignorespaces ``A sobering thought, Eileen: What if, right at this very moment I *am* living up to my full potential?'' -- Jane Wagner}}{39}}
% \newlabel{xxx}{{174}{39}}

\begin{figure}[!htb]
This figure is not -- the difference is the \verb+\label+ command
\caption{\label{yyy}``A sobering thought, Eileen: What if, right at this very
moment I *am* living up to my full potential?'' -- Jane Wagner}
\end{figure}

\captionsetup[table]   {position=top,aboveskip=5pt} % did not work in caption 3.0

\DeclareCaptionFormat{normal}  {#1#2#3\par}
\captionsetup{format=normal}

%-----------------------------------------------------------------------------%

\ifundefined{SCfigure}
---
\else

% Minimalbeispiel fuer folgendes Problem:
% Sobald die Option 'flushleft' statt 'normal' in caption2
% gewaehlt wird, wird nicht mehr die Unterkante der Legende
% mit dem Bild auf eine Linie gesetzt.
% Markus Wellen - markus.wellen@tu-clausthal.de
%
%\captionsetup{style=default,position=default,aboveskip=10pt,belowskip=0pt}% alles in Ordnung
\captionsetup{justification=raggedright,margin=10pt}% gewuenscht, geht aber schief
%\usepackage[rightcaption]{sidecap}
\setlength\unitlength{1cm}
\newcommand{\FIGi}{%
   \fbox{%
     \begin{picture}(4,4)
       \put(2,3){\circle{1}}
       \put(1,1){\circle{1}}
       \put(3,1){\circle{1}}
     \end{picture}}
}
\subsection*{Minimalbeispiel f\"ur ein Problem mit sidecap/caption2}
\begin{SCfigure}[][!htb]
  \FIGi
% \onelinecaptionsfalse
  \caption*{Hier kommt die Legende zum Bild. Wenn sie etwas lang ist und
zudem sehrlangeW\"orter auftauchen, w\"are es besser, Flattersatz gebrauchen
zu k\"onnen.}
\end{SCfigure}
\begin{figure}[!htb]
  \FIGi \FIGi
  \caption{Ein sehr breites Bild. Hier macht es keinen Sinn, mit sidecap
zu arbeiten. %
    Alles, was zur Legende geh\"ort, steht unter dem Bild.
UmdasFlushleftzeigenzuk\"onnen, kommt hier ein ganz langes Wort.}
\end{figure}

\fi

%-----------------------------------------------------------------------------%

\captionsetup{style=default,position=default,aboveskip=10pt,belowskip=0pt}


%************************************%
%*                                  *%
%* Known bugs                       *%
%*                                  *%
%************************************%

\clearpage
\section{Known bugs}

This should produce a entry in the List Of Tables:\\
|\captionof{table}{}|

\begingroup
\captionof{table}{} % BUG! PRODUCES NO LOT!
\endgroup

%\captionof{figure}{} % BUG! PRODUCES NO LOT!

\end{document}
\endinput
